% Narrative rewrite candidate. Not included by main.tex until explicitly swapped.
\section{Related Work}
\label{sec:related-v2}

The communication decision in networked coordination has two distinct
questions: \emph{is the current state informative enough to transmit?}, and
\emph{what does the sender causally know about the state currently held by the
receiver?} Event-triggered control mainly addresses the first question,
whereas age-based scheduling and feedback-aware status updating mainly address
the second. Our problem lies at their intersection. We therefore organize the
literature by the information used to make a transmission decision rather
than by application domain alone.

\subsection{State relevance: event-triggered coordination}

Classical event-triggered control replaces periodic updates with transmissions
driven by a local state or estimation error
\cite{astrom2002riemann,tabuada2007event,heemels2012introduction}. Distributed
extensions establish consensus and coordination guarantees, exclude Zeno
behavior, and introduce dynamic thresholds
\cite{dimarogonas2012distributed,seyboth2013eventbased,girard2015dynamic,yi2017distributed,nowzari2019eventtriggered}.
This line has been extended to multi-UAV formation under directed topologies,
delays, disturbances, and dynamic triggers
\cite{yin2023eventbased,ji2023dynamic,chen2024distributed,yang2025fencing}.
Importantly, Wang \emph{et al.} also demonstrate finite-time distributed
event-triggered formation using three Bebop2 quadrotors
\cite{wang2022finiteuav}. Thus, neither event-triggered multi-agent control nor
its application to UAV formation is new.

What this strand generally abstracts is receiver observability. A sender
usually resets its local broadcast error when it transmits, implicitly treating
the transmitted sample as the neighbors' new reference despite queueing,
forward loss, or delayed feedback. Kesper \emph{et al.} explicitly attach age
timers to last-broadcast states in learned distributed event-triggered control
\cite{kesper2023toward}, but the timers reset on broadcast rather than on
delayed acknowledgement of a particular receiver. State relevance therefore
answers when the plant has changed, but not necessarily whether a specific
neighbor has received that change.

\subsection{Time relevance: AoI-aware networked control}

Age of information (AoI) quantifies the time elapsed since generation of the
newest update available at a receiver
\cite{kaul2012realtime,sun2017update,yates2021aoisurvey}. AoI-aware networked
control uses this receiver-side freshness to schedule scarce transmissions or
constrain estimation quality. Mamduhi \emph{et al.} combine age-based and
event-based scheduling across multiple independent feedback loops
\cite{mamduhi2020freshness}; Wang \emph{et al.} impose freshness constraints in
an event-triggered Kalman consensus filter \cite{wang2021freshness}; and Lin
\emph{et al.} integrate AoI into event-triggered load-frequency control
\cite{lin2023eventtriggered}. These studies establish that freshness can be a
control-relevant scheduling variable. They also show why a pure state trigger
can leave a receiver stale even when the sender's instantaneous innovation is
small.

Age alone, however, is not equivalent to task relevance. In networked control,
AoI can underperform policies that account for plant state or value of
information \cite{ayan2019aoivoi,mamduhi2020freshness}. The communication
literature reaches the same conclusion through semantic freshness metrics.
Age of Incorrect Information (AoII) penalizes information that is both stale
and wrong \cite{maatouk2020aoii}, while Age of Changed Information (AoCI)
combines elapsed time with whether the monitored content has changed
\cite{wang2022aoci}. Decentralized AoII policies have also been studied under
slotted-ALOHA collisions \cite{nayak2023decentralizedaoii}. These works are
important conceptual precedents for combining age and innovation; consequently,
our novelty cannot be the generic idea of making freshness content-aware.
Their typical setting is a discrete status source or remote estimator with a
common destination, rather than continuous-state distributed formation with a
separate belief for every directed neighbor link.

\subsection{Receiver observability: ACKs and delayed two-way feedback}

When delivery is unreliable, the sender does not know the receiver's AoI from
its own transmission history alone. ACK/NACK-aware scheduling addresses this
causal information constraint. Ceran \emph{et al.} optimize AoI through
ACK-driven sampling and retransmission decisions under resource constraints
\cite{ceran2019average}. Pan \emph{et al.} derive threshold sampling for an
unreliable forward channel with random two-way delay and acknowledgement
feedback \cite{pan2023twoway}. Tahir \emph{et al.} go further to decentralized
sensor agents sharing capacity-limited, non-FIFO duplex channels: each agent
uses a particle filter to maintain a belief over receiver AoI under delayed
ACKs \cite{tahir2024collaborative}. Onozuka \emph{et al.} provide a related
bidirectional precedent by event-triggering both forward and feedback traffic
in railway control \cite{onozuka2024aoi}.

This strand solves part of the observability problem that classical
event-triggered control omits, and Tahir \emph{et al.} is therefore a necessary
nearest-neighbor baseline rather than merely a citation. Nevertheless, these
policies optimize age or estimation objectives for sources communicating with
a receiver; they do not decide whether a continuous physical state is newly
informative for a particular formation neighbor. Conversely, AoI-aware control
studies that assume the receiver age is directly available do not expose the
distinction between a state that was transmitted and a state whose reception
has been causally confirmed.

\subsection{UAV shared-medium implementations}

UAV-network studies emphasize mobility, interference, bandwidth constraints,
and topology changes \cite{gupta2016survey,zeng2016wireless,mozaffari2019tutorial}.
WiSwarm is especially relevant because it implements an AoI-based application
layer over Wi-Fi for collaborative UAV teams, including collision prevention,
stale-packet handling, traffic prioritization, and hardware flight experiments
\cite{tripathi2023wiswarm}. It demonstrates that a useful freshness policy
cannot be evaluated only by counting logical sender--receiver updates: shared
medium access and queueing can change both delivered freshness and the relative
cost of broadcast and unicast traffic.

The UAV literature therefore exposes a complementary split. Formation-control
papers provide stronger stability and plant results but often idealize the
network, whereas network-oriented systems such as WiSwarm model or measure the
wireless stack but do not implement the present per-neighbor, state-innovation
formation trigger. Our current simulation addresses the protocol logic between
these layers, but its broadcast metric remains an accounting proxy; without a
contention/MAC or trace-based experiment, it cannot support a shared-medium
performance claim.

\subsection{Unresolved intersection and positioning}

The literature has therefore solved four neighboring subproblems: state-error
triggers for distributed coordination; receiver-age scheduling for networked
control; semantic freshness that combines time and content; and causal
receiver-age inference from delayed ACKs. We did not identify prior work that
jointly combines (i) a per-directed-neighbor estimate updated only by causal
cumulative ACKs, (ii) state innovation measured against the last transmitted
state, (iii) an innovation threshold modulated by the ACK-derived receiver-age
estimate, and (iv) protocol logic that separates genuinely new information
from refresh traffic so that in-flight suppression applies only to the latter,
within a distributed formation-control loop. This is a scoped intersection
claim, not a claim that any individual ingredient is new.

That gap motivates three contributions. First, we formulate a dual-memory
sender policy that distinguishes last-transmitted from ACK-confirmed state and
therefore remains causal under forward and reverse delay/loss. Second, we
separate new-information and refresh branches; a frozen design ablation tests
the failure caused by applying one cooldown to both. Third, we evaluate the
result under matched network realizations and multiple communication-accounting
models, retaining the regimes in which periodic baselines dominate. These
contributions do not constitute a stability theorem, a MAC-layer solution, or
a hardware validation; those boundaries delimit the present study and define
the next validation stage.

